\documentclass[12pt, a4paper]{article}

% Encoding and Language
\usepackage[utf8]{inputenc}
\usepackage[T1]{fontenc}
\usepackage[english]{babel}

% Page Layout
\usepackage[a4paper, margin=2.5cm]{geometry}

% Typography
\usepackage{microtype}
\usepackage[parfill]{parskip}

% Math
\usepackage{amsmath}
\usepackage{amssymb}

% Graphics and Floats
\usepackage{graphicx}
\usepackage{float}

% Tables
\usepackage{booktabs}

% Lists
\usepackage{enumitem}

% Colors
\usepackage[dvipsnames]{xcolor}

% Code Snippets
\usepackage{listings}
\definecolor{codegreen}{rgb}{0,0.6,0}
\definecolor{codegray}{rgb}{0.5,0.5,0.5}
\definecolor{codepurple}{rgb}{0.58,0,0.82}
\definecolor{backcolour}{rgb}{0.95,0.95,0.92}
\lstdefinestyle{mystyle}{
    backgroundcolor=\color{backcolour},
    commentstyle=\color{codegreen},
    keywordstyle=\color{magenta},
    numberstyle=\tiny\color{codegray},
    stringstyle=\color{codepurple},
    basicstyle=\ttfamily\footnotesize,
    breakatwhitespace=false, breaklines=true, captionpos=b,
    keepspaces=true, numbers=left, numbersep=5pt,
    showspaces=false, showstringspaces=false, showtabs=false, tabsize=2
}
\lstset{style=mystyle}

% Hyperlinks and PDF Metadata
\usepackage{hyperref}
\hypersetup{
    colorlinks = true,
    linkcolor  = MidnightBlue,
    citecolor  = MidnightBlue,
    urlcolor   = MidnightBlue,
    pdftitle   = {Practical Report 3: Talk Allocation with CHC Metaheuristic},
    pdfauthor  = {Jesús Fernández López and Rafael Carlos Schoettel Vilchez},
    pdfsubject = {Metaheuristics},
    pdfkeywords= {CHC metaheuristic HUX crossover talk allocation school assignment},
}

% Global Settings
\widowpenalty=10000
\clubpenalty=10000
\frenchspacing

\newcommand{\reporttitle}{Practical Assignment 3: Talk Allocation\\[0.3cm]
\large using a CHC Metaheuristic}
\newcommand{\reportdate}{\today}
\newcommand{\authorA}{Jesús Fernández López}
\newcommand{\authorB}{Rafael Carlos Schoettel Vilchez}

\begin{document}

\begin{titlepage}
  \centering
  \vspace*{3cm}
  {\Huge \textbf{Practical Assignment Report}}\\[0.8cm]
  {\LARGE \textbf{\reporttitle}}\\[0.5cm]
  \rule{\textwidth}{0.4pt}\\[2cm]
  {\Large \textbf{Authors}}\\[0.4cm]
  {\large \authorA\ \& \authorB}\\[1.5cm]
  {\Large \textbf{Subject}}\\[0.4cm]
  {\large Metaheuristics}\\[0.4cm]
  {\normalsize University of Córdoba}\\[1.5cm]
  {\Large \textbf{Date}}\\[0.4cm]
  {\large \reportdate}
  \vfill
  \rule{\textwidth}{0.4pt}
\end{titlepage}

\newpage
\tableofcontents
\newpage

% ============================================================
\section{Problem Analysis}
% ============================================================

\subsection{Context and Motivation}

Every year, to celebrate the \textit{International Day of Women and Girls in Science}
(February 11\textsuperscript{th}), the University of Córdoba organises a programme of
scientific outreach talks at local schools. Schools from the city and the surrounding
province submit requests specifying the educational level and preferred topic of the
desired talk. At the same time, researchers volunteer to deliver talks within their
area of expertise, subject to constraints on travel and educational level.

Once both sets of requests are received, a coordinator must match talks to researchers.
The matching must satisfy hard feasibility requirements while optimising a set of
priority-based soft objectives. The number of researchers $R$ and the number of
requested talks $T$ vary each year, producing two qualitatively different regimes that
the algorithm must handle transparently.

\subsection{Formal Problem Definition}

\subsubsection{Sets and Parameters}

\begin{center}
\renewcommand{\arraystretch}{1.3}
\begin{tabular}{cl}
  \toprule
  \textbf{Symbol} & \textbf{Definition} \\
  \midrule
  $T$ & Total number of talks requested by schools \\
  $E$ & Total number of schools \\
  $R$ & Total number of researchers \\
  $\mathcal{T} = \{t_0, \ldots, t_{T-1}\}$ & Set of requested talks \\
  $\mathcal{S} = \{s_1, \ldots, s_E\}$     & Set of schools \\
  $\mathcal{R} = \{r_1, \ldots, r_R\}$     & Set of researchers \\
  \bottomrule
\end{tabular}
\end{center}

\subsubsection{Entity Attributes}

Each \textbf{school} $s \in \mathcal{S}$ is characterised by:
\begin{itemize}[itemsep=0.1em]
  \item $\text{loc}(s) \in \{\text{city}, \text{province}\}$ — geographical location.
  \item $\text{disadv}(s) \in \{0,1\}$ — whether $s$ is in a disadvantaged area.
  \item $\text{type}(s) \in \{\text{public}, \text{concerted}, \text{private}\}$ — institution type.
  \item $\text{new}(s) \in \{0,1\}$ — whether $s$ participates for the first time.
\end{itemize}

Each \textbf{talk request} $t \in \mathcal{T}$ carries:
\begin{itemize}[itemsep=0.1em]
  \item $\text{school}(t) \in \mathcal{S}$ — the requesting school.
  \item $\text{topic}(t) \in \mathcal{K} \cup \{\text{any}\}$ — required topic ($\mathcal{K}$
        is the set of known disciplines; ``any'' means no restriction).
  \item $\text{level}(t) \in \mathcal{L}$ — required educational level
        (preschool, primary, secondary, high school, or vocational training).
\end{itemize}

Each \textbf{researcher} $r \in \mathcal{R}$ is described by:
\begin{itemize}[itemsep=0.1em]
  \item $\text{topic}(r) \in \mathcal{K}$ — area of expertise.
  \item $\text{level}(r) \in \mathcal{L}$ — preferred educational level.
  \item $\text{travel}(r) \in \{0,1\}$ — whether $r$ can travel outside the city.
  \item $\text{newR}(r) \in \{0,1\}$ — first year participating.
  \item $\text{prevProv}(r) \in \{0,1\}$ — gave a talk in the province last year.
  \item $\text{prevSchool}(r) \in \mathcal{S} \cup \{\varnothing\}$ — school visited last year.
  \item $\text{maxTalks}(r) \in \{1, 2\}$ — maximum talks $r$ is willing to give.
\end{itemize}

\subsubsection{Decision Variable}

The problem is equivalent to a \textbf{Generalised Assignment Problem} (GAP). The
decision variable is an assignment function:
\begin{equation}
  x : \mathcal{T} \to \mathcal{R} \cup \{-1\}
\end{equation}
where $x(t) = r$ means talk $t$ is assigned to researcher $r$, and $x(t) = -1$
means talk $t$ is left unassigned.

\subsubsection{Hard Constraints}

\begin{enumerate}[label=\textbf{HC\arabic*.}, leftmargin=3em, itemsep=0.3em]
  \item \textbf{Topic compatibility:} $\text{topic}(t) = \text{any}$ \,\textit{or}\,
        $\text{topic}(r) = \text{topic}(t)$, for all assigned $(t, r)$.
  \item \textbf{Level compatibility:} $\text{level}(r) = \text{level}(t)$, for all
        assigned $(t, r)$.
  \item \textbf{Coverage when $R \geq T$:} Every school $s$ must have at least one
        assigned talk: $\exists\, t\!:\,\text{school}(t)=s,\;x(t)\neq -1$.
  \item \textbf{Location feasibility:} If $\text{loc}(\text{school}(t)) = \text{province}$
        then $\text{travel}(x(t)) = 1$.
  \item \textbf{Capacity:} $|\{t : x(t) = r\}| \leq \text{maxTalks}(r)$ for all $r$.
\end{enumerate}

\subsubsection{Soft Constraints (Priorities)}

When $T > R$ (insufficient researchers), schools are prioritised in the following order:
\begin{enumerate}[label=\arabic*., itemsep=0.1em]
  \item School participates for the first time.
  \item School is located in the province.
  \item School is in a disadvantaged area.
  \item School is a public institution.
  \item School is concerted.
  \item School is private (lowest priority).
\end{enumerate}

When $R > T$ (surplus researchers), researchers are prioritised:
\begin{enumerate}[label=\arabic*., itemsep=0.1em]
  \item Researcher participates for the first time.
  \item Researcher can travel to the province.
\end{enumerate}

Additional historical soft constraints:
\begin{itemize}[itemsep=0.1em]
  \item Researcher who was in the province last year should preferably go to the city.
  \item Researcher should preferably not revisit the same school as last year.
\end{itemize}

% ============================================================
\section{Design Decisions}
% ============================================================

\subsection{Chromosome Representation}

We encode a solution as a one-dimensional integer array of length $T$:
\begin{equation}
  \mathbf{c} = [c_0,\, c_1,\, \ldots,\, c_{T-1}], \quad c_i \in \{-1,\, 0,\, \ldots,\, R-1\}
\end{equation}

Position $i$ of the chromosome corresponds to talk $t_i$. The value $c_i$ is the
\emph{index} of the researcher assigned to that talk (or $-1$ for unassigned). This
\textbf{direct encoding} has several desirable properties:
\begin{itemize}[itemsep=0.1em]
  \item Compact representation: $T$ genes, each with at most $R+1$ possible values.
  \item Simple crossover: any valid gene-level operator is applicable.
  \item Transparent decoding: solution quality is immediately readable from the array.
\end{itemize}

\subsection{Preprocessing: Valid Researchers per Talk}

Before launching the metaheuristic, we pre-compute for each talk $t$ the set
$V(t)$ of researchers that can feasibly cover it (stored as a Python dictionary
called \texttt{valid\_r\-es\-ear\-chers\_per\_talk}):
\begin{align}
  V(t) = \{\, r \in \mathcal{R} \mid\;
            &(\text{topic}(t) = \text{any} \lor \text{topic}(r) = \text{topic}(t))\nonumber\\
            &\land\; \text{level}(r) = \text{level}(t)\nonumber\\
            &\land\; (\text{loc}(\text{school}(t)) = \text{city} \lor \text{travel}(r) = 1)
            \,\}\nonumber
\end{align}

This \textbf{mandatory preprocessing step}:
\begin{itemize}[itemsep=0.1em]
  \item Enforces hard constraints HC1–HC4 \emph{before} the search begins, drastically
        reducing the effective search space.
  \item Speeds up initialisation: random chromosomes are built by sampling from $V(t)$,
        producing feasible or near-feasible individuals from generation 0.
  \item Makes the fitness function lighter: constraint HC1–HC4 violations can be detected
        in $O(1)$ using the precomputed set.
\end{itemize}

Talks for which $V(t) = \varnothing$ are inherently infeasible; such genes are set to
$-1$ and are never modified by crossover or mutation.

\begin{lstlisting}[language=Python, caption={Preprocessing: building the valid researchers dictionary}]
def build_valid_researchers_per_talk(talks, researchers, schools):
    valid = {}
    for talk in talks:
        school = schools[talk.school_id]
        candidates = []
        for r in researchers.values():
            # Hard constraint 1: topic match
            if talk.topic != "any" and r.topic != talk.topic:
                continue
            # Hard constraint 2: level match
            if r.level != talk.level:
                continue
            # Hard constraint 3: travel feasibility
            if school.location == "province" and not r.can_travel:
                continue
            candidates.append(r.researcher_id)
        valid[talk.talk_id] = candidates
    return valid
\end{lstlisting}

% ============================================================
\section{Fitness Function}
% ============================================================

\subsection{Minimisation Approach}

The fitness function measures the \emph{total penalty} of a chromosome. Following a
minimisation paradigm, a perfect solution has $f(\mathbf{c}) = 0$, and each
unsatisfied constraint contributes a positive penalty proportional to its weight:

\begin{equation}
  f(\mathbf{c}) = \underbrace{\sum_{k} w_k^{\text{hard}} \cdot \delta_k^{\text{hard}}(\mathbf{c})}_{\text{hard penalties}}
  \;+\;
  \underbrace{\sum_{k} w_k^{\text{soft}} \cdot \delta_k^{\text{soft}}(\mathbf{c})}_{\text{soft penalties}}
\end{equation}

where $\delta_k(\mathbf{c}) \in \{0,1,\mathbb{Z}^+\}$ counts violations of constraint $k$.

\subsection{Modular Penalty Architecture}

Each constraint class has its own function, collecting penalties independently and
summing them. This modularity simplifies testing, enables weight tuning, and allows
constraints to be enabled or disabled individually.

\subsubsection{Hard Constraint Penalties}

\renewcommand{\arraystretch}{1.3}
\begin{table}[H]
  \centering
  \caption{Hard constraint penalties and default weights.}
  \label{tab:hard_penalties}
  \resizebox{\textwidth}{!}{%
  \begin{tabular}{llc}
    \toprule
    \textbf{Violation} & \textbf{Description} & \textbf{Default Weight} \\
    \midrule
    School without talk ($R \geq T$) & A school receives zero talks when there are enough researchers & $1{,}000{,}000$ \\
    Location mismatch & Researcher without travel clearance assigned to a province school & $1{,}000{,}000$ \\
    Overallocation & Researcher assigned more talks than their \texttt{maxTalks} limit & $1{,}000{,}000$ \\
    \bottomrule
  \end{tabular}}
\end{table}
\renewcommand{\arraystretch}{1}

Setting hard penalties to $10^6$ ensures that any feasible solution — even one that
satisfies no soft constraint — outranks any infeasible solution.

\subsubsection{Soft Constraint Penalties}

\renewcommand{\arraystretch}{1.3}
\begin{table}[H]
  \centering
  \caption{Soft constraint penalties, their context, and default weights.}
  \label{tab:soft_penalties}
  \resizebox{\textwidth}{!}{%
  \begin{tabular}{llc}
    \toprule
    \textbf{Violation} & \textbf{Description} & \textbf{Default Weight} \\
    \midrule
    First-year school unserved    & School in its first year receives no talk          & 500 \\
    Province school unserved      & Province school receives no talk ($T > R$)         & 300 \\
    Disadvantaged area unserved   & Disadvantaged school receives no talk              & 400 \\
    Public school unserved        & Public institution receives no talk                & 200 \\
    Concerted school unserved     & Concerted school receives no talk                  & 100 \\
    First-year researcher unused  & New researcher not selected ($R > T$)              & 200 \\
    Traveller unused in province  & Travelling researcher not used for province ($R > T$) & 150 \\
    Same school as last year      & Researcher revisits their school from last year    & 300 \\
    Same province as last year    & Province researcher goes to province again         & 150 \\
    \bottomrule
  \end{tabular}}
\end{table}
\renewcommand{\arraystretch}{1}

\subsection{Configuration Dictionary}

All weights are collected in a single Python dict (\texttt{DEFAULT\_CONFIG}).
This makes the fitness function \emph{fully configurable}: a user can pass a custom
dict to \texttt{compute\_fitness()} to reflect different institutional priorities without
modifying any source code.

\begin{lstlisting}[language=Python, caption={Configurable weight dictionary}]
DEFAULT_CONFIG = {
    # Hard constraints
    "w_school_no_talk":     1_000_000,
    "w_location_mismatch":  1_000_000,
    "w_overallocation":     1_000_000,
    # Soft: unserved schools
    "w_first_year_school":  500,
    "w_province_school":    300,
    "w_disadvantaged":      400,
    "w_public_institution": 200,
    "w_concerted":          100,
    # Soft: unused researchers (R > T regime)
    "w_first_year_researcher": 200,
    "w_travel_researcher":     150,
    # Soft: historical
    "w_same_school_as_last_year":   300,
    "w_same_province_as_last_year": 150,
}
\end{lstlisting}

% ============================================================
\section{Metaheuristic Logic}
% ============================================================

\subsection{Overview of CHC}

CHC (Cross-generational elitist selection, Heterogeneous recombination, Cataclysmic
mutation) is an evolutionary algorithm proposed by Eshelman (1991) that combines
aggressive elitism with operators designed to maintain population diversity. Its
three defining mechanisms are:
\begin{enumerate}[itemsep=0.2em]
  \item \textbf{Cross-generational elitism (C):} All individuals — old and new — compete
        for survival; only the best $N_p$ are kept.
  \item \textbf{HUX crossover (H):} A specialised operator that maximises diversity in
        offspring while preserving strong genes from both parents.
  \item \textbf{Cataclysmic restart (C):} When the population has converged completely,
        a controlled restart diversifies around the best individual.
\end{enumerate}

Unlike standard GAs, CHC uses \emph{no mutation during the main loop}. Instead,
incest prevention prevents premature convergence, and the restart provides macro-level
diversity when the incest threshold reaches zero.

\subsection{HUX Crossover}

The Half-Uniform Crossover operator is the core recombination method in CHC. Given
two parents $\mathbf{p}_1$ and $\mathbf{p}_2$, HUX:
\begin{enumerate}[itemsep=0.1em]
  \item Identifies the set $D$ of positions where $p_1$ and $p_2$ differ:
        $D = \{i : c_i^{(1)} \neq c_i^{(2)}\}$.
  \item Randomly selects exactly $\lfloor |D|/2 \rfloor$ positions from $D$ to swap.
  \item Produces two offspring $\mathbf{c}_1$, $\mathbf{c}_2$ by exchanging the
        selected genes between parents.
\end{enumerate}

\begin{equation}
  |D(\mathbf{c}_1, \mathbf{c}_2)| = |D(\mathbf{c}_1, \mathbf{p}_1)|
  = |D(\mathbf{c}_2, \mathbf{p}_2)|
  = \left\lfloor \frac{|D(\mathbf{p}_1, \mathbf{p}_2)|}{2} \right\rfloor
\end{equation}

This property guarantees that each offspring is exactly halfway between the two parents
in Hamming space, which is the furthest apart they can be while each carrying a mix of
both parents. HUX maximises exploration among the offspring of a given pair.

\begin{lstlisting}[language=Python, caption={HUX crossover implementation}]
def hux_crossover(parent1, parent2):
    diff_positions = [i for i, (a, b) in enumerate(zip(parent1, parent2)) if a != b]
    if len(diff_positions) < 2:
        return copy.copy(parent1), copy.copy(parent2)
    half = len(diff_positions) // 2
    swap_positions = set(random.sample(diff_positions, half))
    child1, child2 = copy.copy(parent1), copy.copy(parent2)
    for pos in swap_positions:
        child1[pos], child2[pos] = child2[pos], child1[pos]
    return child1, child2
\end{lstlisting}

\subsection{Incest Prevention and the Hamming Distance Threshold}

A key feature of CHC is the \textbf{incest prevention} mechanism. Two individuals are
only allowed to mate if their Hamming distance exceeds a dynamic threshold $d$:

\begin{equation}
  d_H(\mathbf{p}_1, \mathbf{p}_2) = \sum_{i=0}^{T-1} \mathbf{1}[c_i^{(1)} \neq c_i^{(2)}] > d
\end{equation}

The threshold $d$ is initialised to $\lfloor T/4 \rfloor$ and decreases by 1 each
generation in which no offspring improves over both parents. This progressive relaxation
allows the algorithm to increase mating pressure as the population converges, eventually
reaching all pairs.

\textbf{Elitist replacement:} Children only enter the population if they strictly
improve over at least one of their parents. This cross-generational competition ensures
monotonic improvement: the best fitness never decreases.

\subsection{Cataclysmic Restart}

When $d$ reaches $0$ and no crossover produces improvements, the population has
converged. Rather than terminating, CHC performs a \textbf{cataclysmic restart}:

\begin{enumerate}[itemsep=0.1em]
  \item The current best individual (the \emph{survivor}) is copied unchanged.
  \item The remaining $N_p - 1$ individuals are created by randomly mutating exactly
        $\lfloor \mu \cdot T \rfloor$ genes of the survivor ($\mu = 35\%$ by default).
  \item Each mutated gene is replaced by a valid researcher drawn uniformly from
        $V(t)$ for that position.
  \item The Hamming threshold is reset to $\lfloor T/4 \rfloor$.
\end{enumerate}

This restart is \emph{controlled}: the best solution is never lost, but the population
is aggressively diversified around it, enabling further exploration in new regions.

\begin{lstlisting}[language=Python, caption={Cataclysmic restart}]
def cataclysmic_mutation(best, pop_size, mutation_rate, talks, valid_map, r_id_list):
    new_population = [copy.copy(best)]
    n_mutate = max(1, int(mutation_rate * len(talks)))
    for _ in range(pop_size - 1):
        mutant = copy.copy(best)
        positions = random.sample(range(len(talks)), n_mutate)
        for pos in positions:
            mutant[pos] = repair_gene(pos, valid_map, r_id_list)
        new_population.append(mutant)
    return new_population
\end{lstlisting}

\subsection{Algorithm Flow}

Algorithm~\ref{alg:chc} summarises the complete CHC loop.

\begin{figure}[H]
\centering
\begin{minipage}{0.92\textwidth}
\begin{lstlisting}[language=Python, caption={CHC main loop (simplified)}, label={alg:chc}]
threshold = T // 4
population = initialise_population(pop_size)
while generation < max_generations:
    shuffle(population)
    pairs = [(pop[i], pop[i+1]) for i in range(0, pop_size-1, 2)]
    new_individuals = []
    for p1, p2 in pairs:
        if hamming(p1, p2) > threshold:
            c1, c2 = hux_crossover(p1, p2)
            if fitness(c1) < min(fitness(p1), fitness(p2)):
                new_individuals.append(c1)
            if fitness(c2) < min(fitness(p1), fitness(p2)):
                new_individuals.append(c2)
    if new_individuals:
        merge and keep top pop_size individuals   # cross-generational elitism
    else:
        threshold -= 1                            # tighten incest threshold
    if threshold <= 0:
        population = cataclysmic_restart(best)   # reset around best
        threshold = T // 4
    generation += 1
\end{lstlisting}
\end{minipage}
\end{figure}

\subsection{Parameter Configuration}

\renewcommand{\arraystretch}{1.3}
\begin{table}[H]
  \centering
  \caption{CHC hyperparameters used in this implementation.}
  \label{tab:params}
  \begin{tabular}{lcc}
    \toprule
    \textbf{Parameter} & \textbf{Symbol} & \textbf{Default Value} \\
    \midrule
    Population size        & $N_p$  & 50 \\
    Maximum generations    & $G$    & 200 \\
    Initial threshold      & $d_0$  & $\lfloor T/4 \rfloor$ \\
    Restart mutation rate  & $\mu$  & 35\% \\
    \bottomrule
  \end{tabular}
\end{table}
\renewcommand{\arraystretch}{1}

The choice of $\mu = 35\%$ follows Eshelman's recommendation: too small a mutation
produces restarts that are virtually identical to the best, while too large a value
effectively discards the best solution's genetic information.

% ============================================================
\section{Conclusions}
% ============================================================

This practice addressed the \textbf{Talk Allocation Problem} as a Generalised
Assignment Problem with one hard categorical constraint per dimension (topic,
level, location) and a rich set of lexicographic soft priorities. The CHC
metaheuristic was selected because it naturally balances intense elitism with
controlled diversity preservation, without relying on hand-tuned mutation
probabilities during the main evolutionary loop.

The modular penalty architecture allows practitioners at the University of Córdoba
to adjust the relative importance of each priority rule — for example, giving more
weight to disadvantaged areas in a particular year — without touching the algorithm's
source code.

The mandatory preprocessing phase (\texttt{valid\_researchers\_per\_talk}) reduces
feasibility enforcement to an $O(1)$ lookup during fitness evaluation, making the
algorithm scalable to larger real-world instances with hundreds of talks and researchers.

Future work could introduce a \emph{repair operator} that, whenever a child violates a
hard constraint on a gene, replaces that gene with a draw from $V(t)$ before
evaluating fitness. This would further accelerate convergence by keeping the
population closer to the feasible region at all times.

\end{document}
